\documentclass[a4paper,11pt]{ctexart}
\usepackage[top=2cm, bottom=2cm, left=2cm, right=2cm]{geometry}
\usepackage{titlesec}
\usepackage{enumitem}
\usepackage{multicol}
\usepackage{fancyhdr}
\usepackage{tcolorbox}

% 页面设置
\pagestyle{fancy}
\fancyhf{}
\fancyhead[C]{思政课复习资料汇总 (第10-16讲)}
\fancyfoot[C]{\thepage}

% 标题格式设置
\titleformat{\section}{\Large\bfseries\centering}{}{0em}{}
\titleformat{\subsection}{\large\bfseries}{\thesubsection}{1em}{}

% 列表设置
\setlist[enumerate,1]{label=\arabic*., itemsep=2pt, parsep=0pt}
\setlist[enumerate,2]{label=\Alph*., itemsep=0pt, parsep=0pt}

% 自定义答案框
\newtcolorbox{answerbox}{colback=gray!10!white, colframe=gray!50!black, title=参考答案}

\title{\textbf{《习近平新时代中国特色社会主义思想概论》\\期末复习资料汇总}}
\author{整理}
\date{\today}

\begin{document}

\maketitle


% ==========================================================================================
% 第十讲
% ==========================================================================================
\section{第十讲 加强以民生为重点的社会建设 [cite: 1]}

\subsection{重要知识点}
\begin{enumerate}
    \item 民生是最大的政治
    \item 增进民生福祉是发展的根本目的
    \item 脱贫攻坚精神
    \item 坚持系统治理、综合治理、依法治理、源头治理
    \item 健全党组织领导的自治、法治、德治相结合的城乡基层治理体系
\end{enumerate}

\subsection{练习题}
\textbf{（一）单选题} \\
\textit{（注：原始资料未提供题目文本，仅提供参考答案）}

\begin{answerbox}
1.C \quad 2.A \quad 3.B \quad 4.A \quad 5.D \quad 6.A \quad 7.D \quad 8.A \quad 9.C \quad 10.B
\end{answerbox}

\newpage

% ==========================================================================================
% 第十一讲
% ==========================================================================================
\section{第十一讲 建设社会主义生态文明 [cite: 7, 8]}

\subsection{重要知识点}
\begin{enumerate}
    \item 习近平生态文明思想
    \item 生态文明的由来
    \item 人与自然和谐共生的现代化
    \item 绿水青山就是金山银山
    \item 用最严格制度最严密法治保护生态环境
    \item 扎实推进碳达峰碳中和工作
\end{enumerate}

\subsection{一、选择题}

\textbf{（一）单选题}
\begin{enumerate}
    \item 要树立尊重自然、顺应自然、保护自然的生态文明理念，增强（ ）的意识。
    \begin{multicols}{2}
    \begin{enumerate}
        \item 金山银山不如绿水青山
        \item 绿水青山就是金山银山
        \item 金山银山就是绿水青山
        \item 绿水青山胜过金山银山
    \end{enumerate}
    \end{multicols}
    
    \item 加快构建生态文明体系，确保到（ ）美丽中国目标基本实现。
    \begin{multicols}{4}
    \begin{enumerate}
        \item 2020年
        \item 2025年
        \item 2030年
        \item 2035年
    \end{enumerate}
    \end{multicols}
\end{enumerate}

\begin{answerbox}
单选题参考答案：1.B \quad 2.D \quad 3.B \quad 4.A \quad 5.C
\end{answerbox}

\textbf{（二）多选题}
\begin{enumerate}
    \item 在整个发展过程中，我们都要坚持（）的方针，要像保护眼睛一样保护生态环境
    \begin{multicols}{2}
    \begin{enumerate}
        \item 节约优先
        \item 发展为主
        \item 保护优先
        \item 自然恢复为主
    \end{enumerate}
    \end{multicols}

    \item 绿水青山就是金山银山，阐述了（）的关系
    \begin{multicols}{2}
    \begin{enumerate}
        \item 社会和谐
        \item 经济发展
        \item 生态环境保护
        \item 人类生存
    \end{enumerate}
    \end{multicols}

    \item 建设生态文明是关系（ ）、关乎（ ）的千年大计
    \begin{multicols}{2}
    \begin{enumerate}
        \item 人民福祉
        \item 国家发展
        \item 中华民族永续发展
        \item 社会长治久安
    \end{enumerate}
    \end{multicols}

    \item 加快形成绿色生活方式，增强全面（），培养生态道德和行为习惯
    \begin{multicols}{4}
    \begin{enumerate}
        \item 节约意识
        \item 环保意识
        \item 发展意识
        \item 生态意识
    \end{enumerate}
    \end{multicols}

    \item 人类经历了（ ），是工业文明发展到一定阶段的产物，是实现人与自然和谐发展的新要求。
    \begin{multicols}{4}
    \begin{enumerate}
        \item 原始文明
        \item 农业文明
        \item 工业文明
        \item 生态文明
    \end{enumerate}
    \end{multicols}

    \item 实现（ ）是我国向世界作出的庄严承诺，也是一场广泛而深刻的经济社会变革，绝不是轻轻松松就能实现的。
    \begin{multicols}{2}
    \begin{enumerate}
        \item 碳达峰
        \item 碳中和
        \item 建设美丽世界
        \item 构建人与自然的生命共同体
    \end{enumerate}
    \end{multicols}

    \item 推进生态文明建设，解决（ ）的问题，必须采取一些硬措施，真抓实干才能见效。
    \begin{multicols}{2}
    \begin{enumerate}
        \item 资源约束趋紧
        \item 经济社会发展
        \item 环境污染严重
        \item 生态系统退化
    \end{enumerate}
    \end{multicols}

    \item 要加快推进（ ），开创美丽中国建设新局面。
    \begin{multicols}{2}
    \begin{enumerate}
        \item 绿色发展
        \item 循环发展
        \item 低碳发展
        \item 可持续发展
    \end{enumerate}
    \end{multicols}

    \item 党的二十大报告指出：“（ ），是全面建设社会主义现代化国家的内在要求。”
    \begin{multicols}{4}
    \begin{enumerate}
        \item 敬畏自然
        \item 尊重自然
        \item 顺应自然
        \item 保护自然
    \end{enumerate}
    \end{multicols}

    \item 坚持（ ），持续深入打好蓝天、碧水、净土保卫战。
    \begin{multicols}{4}
    \begin{enumerate}
        \item 精准治污
        \item 科学治污
        \item 依法治污
        \item 全民治污
    \end{enumerate}
    \end{multicols}
\end{enumerate}

\subsection{二、简答题}
\textbf{1. 习近平生态文明思想主要包括哪些内容？}
\begin{itemize}
    \item 第一，坚持人与自然和谐共生。
    \item 第二，绿水青山就是金山银山。
    \item 第三，良好生态环境是最普惠的民生福祉。
    \item 第四，统筹山水林田湖草沙系统治理。
    \item 第五，用最严格制度最严密法治保护生态环境。
    \item 第六，共谋全球生态文明建设。
\end{itemize}

\newpage

% ==========================================================================================
% 第十二讲
% ==========================================================================================
\section{第十二讲 建设巩固国防和强大人民军队 [cite: 4, 5]}

\subsection{重要知识点}
\begin{enumerate}
    \item 国际战略形势发生深刻变化
    \item 新时代的强军目标
    \item 国防和军队改革的“三大战役”
    \item 国防和军队改革取得重大成效
    \item 新时代人民军队使命任务
\end{enumerate}

\subsection{一、选择题}

\textbf{（一）单选题}
\begin{enumerate}
    \item 把（ ）作为唯一的根本的标准，是有效履行我军根本职能的内在要求，也是提高军队建设质量效益的客观需要。
    \begin{multicols}{4}
    \begin{enumerate}
        \item 政治建军
        \item 军事改革
        \item 战斗力
        \item 军事创新
    \end{enumerate}
    \end{multicols}
    
    \item （ ）是国家安全的准则。
    \begin{multicols}{4}
    \begin{enumerate}
        \item 政治安全
        \item 国家利益至上
        \item 人民安全
        \item 国际安全
    \end{enumerate}
    \end{multicols}

    \item 党的十九大报告指出，必须全面贯彻党领导人民军队的一系列根本原则和制度，确立 （ ）在国防和军队建设中的指导地位。
    \begin{multicols}{2}
    \begin{enumerate}
        \item 新时代党的建军思想
        \item 新时期党的建军思想
        \item 新时代党的强军思想
        \item 新时期党的强军思想
    \end{enumerate}
    \end{multicols}

    \item 党的十九大报告中指出，组建（ ）机构，维护军人军属合法权益，让军人成为全社会尊崇的职业。
    \begin{multicols}{2}
    \begin{enumerate}
        \item 退伍军人管理保障
        \item 退伍军人社会保障
        \item 退役军人管理保障
        \item 退役军人社会保障
    \end{enumerate}
    \end{multicols}

    \item 力争到二〇三五年\_\_\_\_\_\_国防和军队现代化，到本世纪中叶把人民军队\_\_\_\_\_\_世界一流军队。（ ）
    \begin{multicols}{2}
    \begin{enumerate}
        \item 全面实现；基本建成
        \item 全面实现；全面建成
        \item 基本实现；基本建成
        \item 基本实现；全面建成
    \end{enumerate}
    \end{multicols}

    \item 国家安全是安邦定国的重要基石，（ ）是全国各族人民根本利益所在。
    \begin{multicols}{2}
    \begin{enumerate}
        \item 加快经济发展
        \item 维护国家统一
        \item 促进国际合作
        \item 维护国家安全
    \end{enumerate}
    \end{multicols}

    \item 军队是要准备打仗的，一切工作都必须坚持（ ）标准，向能打仗、打胜仗聚焦。
    \begin{multicols}{4}
    \begin{enumerate}
        \item 战斗力
        \item 斗争力
        \item 战争力
        \item 硬实力
    \end{enumerate}
    \end{multicols}

    \item （ ）是党指挥枪原则落地生根的坚实基础。
    \begin{multicols}{4}
    \begin{enumerate}
        \item 支部建在排上
        \item 支部建在连上
        \item 支部建在营上
        \item 支部建在团上
    \end{enumerate}
    \end{multicols}

    \item 国防和军队改革取得历史性突破，形成军委管总、战区主战、（ ）新格局，人民军队组织架构和力量体系实现革命性重塑。
    \begin{multicols}{4}
    \begin{enumerate}
        \item 党委领导
        \item 战区主建
        \item 军种主战
        \item 军种主建
    \end{enumerate}
    \end{multicols}

    \item 2014年10月，习近平强调革命的（ ）是革命军队的生命线，明确提出军队政治工作的时代主题。
    \begin{multicols}{4}
    \begin{enumerate}
        \item 保障工作
        \item 思想工作
        \item 宣传工作
        \item 政治工作
    \end{enumerate}
    \end{multicols}

    \item （ ）,是决定军队建设的政治方向。
    \begin{multicols}{2}
    \begin{enumerate}
        \item 军民融合是关键
        \item 能打胜仗是核心
        \item 听党指挥是灵魂
        \item 作风优良是保证
    \end{enumerate}
    \end{multicols}

    \item 全党全军全国各族人民要大力弘扬（ ）的光荣传统，不断发展坚如磐石的军政军民关系。
    \begin{multicols}{2}
    \begin{enumerate}
        \item 军爱民
        \item 民拥军
        \item 军爱民、民拥军
        \item 军民团结
    \end{enumerate}
    \end{multicols}

    \item （ ）是军队最重要的战略资源。
    \begin{multicols}{4}
    \begin{enumerate}
        \item 武装装备
        \item 优良传统
        \item 军事人才
        \item 战斗精神
    \end{enumerate}
    \end{multicols}

    \item 2013年11月，党的十八届三中全会把（ ）单列为一个部分写入全会决定中，纳入全面深化改革总体布局、上升为党的意志和国家行为，这在党的历史上是首次。
    \begin{multicols}{2}
    \begin{enumerate}
        \item 作战部署和军队兵种
        \item 深化国防和军队改革
        \item 政治建军和国防建设
        \item 党的领导和军队建设
    \end{enumerate}
    \end{multicols}

    \item 培养“四有”新时代革命军人，锻造“四铁”过硬部队。“四有”指的是（ ）
    \begin{enumerate}
        \item 有理想、有本领、有担当、有希望
        \item 有志气、有底气、有骨气、有勇气
        \item 有灵魂、有本事、有血性、有品德
        \item 有本领、有胆略、有灵魂、有斗志
    \end{enumerate}
\end{enumerate}

\begin{answerbox}
单选题参考答案：
1.C \quad 2.B \quad 3.C \quad 4.C \quad 5.D \quad 6.D \quad 7.A \quad 8.B \quad 9.D \quad 10.D \quad 11.C \quad 12.C \quad 13.C \quad 14.B \quad 15.C
\end{answerbox}

\textbf{（二）多选题}
\begin{enumerate}
    \item 党在新时代的强军目标是建设一支（ ）的人民军队。
    \begin{multicols}{4}
    \begin{enumerate}
        \item 听党指挥
        \item 能打胜仗
        \item 作风优良
        \item 敢于创新
    \end{enumerate}
    \end{multicols}

    \item 推进强军事业必须坚持（ ），更加注重聚焦实战、更加注重创新驱动、更加注重体系建设、更加注重集约高效、更加注重军民融合，全面提高革命化现代化正规化水平。
    \begin{multicols}{4}
    \begin{enumerate}
        \item 政治建军
        \item 改革强军
        \item 科技兴军
        \item 依法治军
    \end{enumerate}
    \end{multicols}

    \item 坚持党对军队绝对领导的主要制度包括（ ）。
    \begin{enumerate}
        \item 军队最高领导权和指挥权属于党中央和中央军委，中央军委实行主席负责制
        \item 实行党委制、政治委员制、政治机关制
        \item 实行党委统一的集体领导下的首长分工负责制
        \item 实行支部建在连上
    \end{enumerate}

    \item 要着力培养新时代革命军人，锻造具有（ ）的过硬部队，确保我军永远立于不败之地。
    \begin{multicols}{2}
    \begin{enumerate}
        \item 铁一般信仰
        \item 铁一般信念
        \item 铁一般纪律
        \item 铁一般担当
    \end{enumerate}
    \end{multicols}

    \item 完善风险防控机制，建立健全（ ），主动加强协调配合，坚持一级抓一级、层层抓落实。
    \begin{multicols}{2}
    \begin{enumerate}
        \item 风险研判机制
        \item 决策风险研判机制
        \item 风险防控协同控制
        \item 风险防控责任机制
    \end{enumerate}
    \end{multicols}

    \item 新时代全军必须夯实（ ）这个强军之基，持之以恒推进作风建设和反腐败斗争，自觉践行人民军队根本宗旨，保持人民军队长期形成的良好形象。
    \begin{multicols}{4}
    \begin{enumerate}
        \item 依法治军
        \item 从严治军
        \item 政治建军
        \item 改革兴军
    \end{enumerate}
    \end{multicols}

    \item 军队要创新发展军事战略指导，构建中国特色现代作战体系，全面提高新时代备战打仗能力有效（ ）
    \begin{multicols}{4}
    \begin{enumerate}
        \item 塑造态势
        \item 管控危机
        \item 遏制战争
        \item 打赢战争
    \end{enumerate}
    \end{multicols}

    \item 必须全面贯彻党领导军队的一系列根本原则和制度，确保军队（ ）。
    \begin{multicols}{4}
    \begin{enumerate}
        \item 绝对忠诚
        \item 绝对纯洁
        \item 绝对可靠
        \item 绝对服从
    \end{enumerate}
    \end{multicols}

    \item 坚持党对军队的绝对领导，就是要确保全军官兵始终在（ ）上与党中央、中央军委保持高度一致，一切行动听从党中央、中央军委指挥。
    \begin{multicols}{4}
    \begin{enumerate}
        \item 政治立场
        \item 政治方向
        \item 政治原则
        \item 政治道路
    \end{enumerate}
    \end{multicols}

    \item 坚持五湖四海、任人唯贤，坚持德才兼备、以德为先，坚持对党忠诚、善谋打仗、敢于担当、实绩突出、清正廉洁的军队好干部标准，完善干部选拔任用机制，增强选人用人的（ ），确保枪杆子始终掌握在忠于党的可靠的人手中。
    \begin{multicols}{4}
    \begin{enumerate}
        \item 科学性
        \item 准确性
        \item 可靠性
        \item 公信度
    \end{enumerate}
    \end{multicols}

    \item 面对强国强军的时代要求，必须切实把战斗力标准在（ ）等各项工作中确立起来。
    \begin{multicols}{4}
    \begin{enumerate}
        \item 军事
        \item 政治
        \item 后勤
        \item 装备
    \end{enumerate}
    \end{multicols}

    \item 武器装备远程（ ）趋势更加明显，战争形态加速向信息化战争演变，智能化战争初现端倪。
    \begin{multicols}{4}
    \begin{enumerate}
        \item 精确化
        \item 智能化
        \item 隐身化
        \item 无人化
    \end{enumerate}
    \end{multicols}

    \item 加强国防教育，增强全民国防观念，使（ ）成为全社会的思想共识和自觉行动。
    \begin{multicols}{4}
    \begin{enumerate}
        \item 关心国防
        \item 热爱国防
        \item 建设国防
        \item 保卫国防
    \end{enumerate}
    \end{multicols}

    \item 为了实施科技兴军战略，必须瞄准世界军事科技前沿，加快（ ）技术发展，争取实现弯道超车，不断提高科技创新对人民军队建设和战斗力发展的贡献率。
    \begin{multicols}{4}
    \begin{enumerate}
        \item 战略性
        \item 前沿性
        \item 颠覆性
        \item 创新性
    \end{enumerate}
    \end{multicols}

    \item 建设世界一流军队，就是要锻造（ ）的精兵劲旅。
    \begin{multicols}{4}
    \begin{enumerate}
        \item 能打硬仗
        \item 召之即来
        \item 来之能战
        \item 战之必胜
    \end{enumerate}
    \end{multicols}

    \item 2014年4月，中央军委印发《关于贯彻落实军委主席负责制建立和完善相关工作机制的意见》，建立（ ）三项机制，推动军委主席负责制各项要求机制化运行。
    \begin{multicols}{4}
    \begin{enumerate}
        \item 行动汇报
        \item 请示报告
        \item 督促检查
        \item 信息服务
    \end{enumerate}
    \end{multicols}

    \item 构建一体化的国家战略体系和能力，是一个系统工程，涉及领域宽、范围广、内容多，必须加强集中统一领导，努力形成（ ）的工作运行体系。
    \begin{multicols}{4}
    \begin{enumerate}
        \item 统一领导
        \item 国家主导
        \item 需求牵引
        \item 市场运作
    \end{enumerate}
    \end{multicols}

    \item 构建一体化的国家战略体系和能力，是一个系统工程，涉及领域宽、范围广、内容多，必须加强集中统一领导，努力形成（ ）的政策制度体系。
    \begin{multicols}{4}
    \begin{enumerate}
        \item 系统完备
        \item 衔接配套
        \item 有效激励
        \item 协调发展
    \end{enumerate}
    \end{multicols}

    \item 党的十八大以来，人民军队重振政治纲纪、（ ），在中国特色强军之路上迈出了坚定步伐。
    \begin{multicols}{2}
    \begin{enumerate}
        \item 重塑组织形态
        \item 重整斗争格局
        \item 重构建设布局
        \item 重树作风形象
    \end{enumerate}
    \end{multicols}

    \item 形成（ ）新格局，打造以精锐作战力量为主体的军事力量体系，初步形成中国特色社会主义军事政策制度体系基本框架。
    \begin{multicols}{4}
    \begin{enumerate}
        \item 军委管总
        \item 统一部署
        \item 战区主战
        \item 军种主建
    \end{enumerate}
    \end{multicols}
\end{enumerate}

\begin{answerbox}
多选题参考答案：
1.ABC \quad 2.ABCD \quad 3.ABCD \quad 4.ABCD \quad 5.ABCD \quad 6.AB \quad 7.ABCD \quad 8.ABC \quad 9.ABCD \quad 10.ABD \quad 11.ABCD \quad 12.ABCD \quad 13.ABCD \quad 14.ABC \quad 15.BCD \quad 16.BCD \quad 17.BCD \quad 18.ABC \quad 19.ABCD \quad 20.ACD
\end{answerbox}

\subsection{二、简答题}
\textbf{1. 如何理解习近平强军思想的主要内容？}
\begin{itemize}
    \item 一是明确强国必须强军，巩固国防和强大人民军队是新时代坚持和发展中国特色社会主义、实现中华民族伟大复兴的战略支撑。
    \item 二是明确党在新时代的强军目标是建设一支听党指挥、能打胜仗、作风优良的人民军队，必须同国家现代化进程相一致，力争到2035年基本实现国防和军队现代化，到本世纪中叶把人民军队全面建成世界一流军队。
    \item 三是明确党对军队的绝对领导是人民军队建军之本、强军之魂，必须全面贯彻党领导军队的一系列根本原则和制度，确保部队绝对忠诚、绝对纯洁、绝对可靠。
\end{itemize}

\newpage

% ==========================================================================================
% 第十三讲
% ==========================================================================================
\section{第十三讲 全面贯彻落实总体国家安全观 [cite: 11, 12]}

\subsection{重要知识点}
\begin{enumerate}
    \item 国家安全
    \item 总体国家安全观
    \item 总体国家安全观的形成过程
    \item 坚持中国特色国家安全道路
    \item 着力防范化解重大风险
\end{enumerate}

\subsection{一、选择题}

\textbf{（一）单选题}
\begin{enumerate}
    \item 坚持总体国家安全观，必须坚持以（ ）为宗旨。
    \begin{multicols}{4}
    \begin{enumerate}
        \item 人民安全
        \item 政治安全
        \item 国内安全
        \item 国际安全
    \end{enumerate}
    \end{multicols}

    \item （ ）是人民幸福安康的基本要求，是安邦定国的重要基石。
    \begin{multicols}{4}
    \begin{enumerate}
        \item 国家富强
        \item 人民富裕
        \item 社会稳定
        \item 国家安全
    \end{enumerate}
    \end{multicols}

    \item 以下不是反映中华优秀传统文化中蕴含着丰富的国家安全战略思想的是（ ）
    \begin{multicols}{2}
    \begin{enumerate}
        \item “生于忧患，死于安乐”
        \item “亲人善邻，国之宝也”
        \item “内事文而和，外事武而义”
        \item “美人之美，美美与共”
    \end{enumerate}
    \end{multicols}

    \item 国家安全工作，归根结底是保障（ ），为群众安居乐业提供坚强保障。有了安全感，获得感才有保障，幸福感才会持久。
    \begin{multicols}{4}
    \begin{enumerate}
        \item 人民利益
        \item 社会稳定
        \item 经济发展
        \item 领土完整
    \end{enumerate}
    \end{multicols}

    \item （ ）设立中央国家安全委员会。
    \begin{multicols}{4}
    \begin{enumerate}
        \item 2013.11
        \item 2014.2
        \item 2016.3
        \item 2016.12
    \end{enumerate}
    \end{multicols}

    \item 每年（ ）被确定为全民国家安全教育日。
    \begin{multicols}{4}
    \begin{enumerate}
        \item 4月15日
        \item 5月17日
        \item 6月15日
        \item 7月30日
    \end{enumerate}
    \end{multicols}

    \item （ ）将“坚持总体国家安全观”纳入新时代坚持和发展中国特色社会主义的基本方略，并写入党章。
    \begin{multicols}{2}
    \begin{enumerate}
        \item 党的十八届三中全会
        \item 党的十九届六中全会
        \item 党的十九大
        \item 党的二十大
    \end{enumerate}
    \end{multicols}

    \item 总体国家安全观强调做好国家安全工作的系统思维和方法，做到（ ），着力解决国家安全工作不平衡不充分的问题。
    \begin{multicols}{4}
    \begin{enumerate}
        \item “四个统筹”
        \item “五个统筹”
        \item “六个统筹”
        \item “七个统筹”
    \end{enumerate}
    \end{multicols}

    \item 健全开放安全保障体系，是“中国开放的大门越开越大”的（ ）。
    \begin{multicols}{4}
    \begin{enumerate}
        \item 重要保证
        \item 基础条件
        \item 根本遵循
        \item 基本原则
    \end{enumerate}
    \end{multicols}

    \item 坚持总体国家安全观，以政治安全为根本。
    \begin{multicols}{4}
    \begin{enumerate}
        \item 人民安全
        \item 政治安全
        \item 国内安全
        \item 国际安全
    \end{enumerate}
    \end{multicols}
    \textit{注：此题题干与选项似乎有出入，原文为“坚持总体国家安全观，以政治安全为根本”，选项为D. 国际安全，请以参考答案D为准或核对题干。}

    \item 走中国特色国家安全道路，努力开创国家安全工作新局面，是对总体国家安全观的贯彻落实，归根到底是为了确保（ ）进程不被迟滞甚至中断。
    \begin{multicols}{2}
    \begin{enumerate}
        \item 社会主义现代化建设
        \item 全体人民共同富裕
        \item 中华民族伟大复兴
        \item 经济社会发展
    \end{enumerate}
    \end{multicols}

    \item （ ）与人民群众切身利益关系最密切，是人民群众安全感的晴雨表，是社会安定的风向标。
    \begin{multicols}{4}
    \begin{enumerate}
        \item 政治安全
        \item 经济安全
        \item 社会安全
        \item 文化安全
    \end{enumerate}
    \end{multicols}

    \item 和平稳定的国际环境和国际秩序是国家安全的（ ）。
    \begin{multicols}{4}
    \begin{enumerate}
        \item 根本方向
        \item 主要目标
        \item 重要保障
        \item 基础条件
    \end{enumerate}
    \end{multicols}

    \item （ ）是中华民族能够生生不息、绵延不绝的文化基因，也是融入中国共产党精神血脉的政治品质。
    \begin{multicols}{4}
    \begin{enumerate}
        \item 以和为贵
        \item 忧患意识
        \item 天下一家
        \item 大同世界
    \end{enumerate}
    \end{multicols}

    \item 面对波谲云诡的国际形势、复杂敏感的周边环境、艰巨繁重的改革发展稳定任务，我们必须树立（ ）。
    \begin{multicols}{4}
    \begin{enumerate}
        \item 忧患意识
        \item 零和思维
        \item 底线思维
        \item 全局意识
    \end{enumerate}
    \end{multicols}
\end{enumerate}

\begin{answerbox}
单选题参考答案：
1.A \quad 2.D \quad 3.D \quad 4.A \quad 5.A \quad 6.A \quad 7.C \quad 8.B \quad 9.A \quad 10.D \quad 11.C \quad 12.C \quad 13.C \quad 14.B \quad 15.C
\end{answerbox}

\textbf{（二）多选题}
\begin{enumerate}
    \item 加强（ ），巩固国家安全人民防线。
    \begin{multicols}{2}
    \begin{enumerate}
        \item 国家安全宣传教育
        \item 全民国防教育
        \item 文化素质教育
        \item 爱国主义教育
    \end{enumerate}
    \end{multicols}

    \item 政治安全的核心是（ ），最根本的就是（ ）。
    \begin{multicols}{2}
    \begin{enumerate}
        \item 政权安全
        \item 制度安全
        \item 维护中国共产党的领导和执政地位
        \item 维护中国特色社会主义制度
    \end{enumerate}
    \end{multicols}

    \item 网络安全已经成为我国面临的（ ）的非传统安全问题之一。
    \begin{multicols}{4}
    \begin{enumerate}
        \item 最普遍
        \item 最复杂
        \item 最现实
        \item 最严峻
    \end{enumerate}
    \end{multicols}

    \item 要坚持（ ）的新安全观，积极塑造外部安全环境
    \begin{multicols}{4}
    \begin{enumerate}
        \item 共同
        \item 综合
        \item 合作
        \item 可持续
    \end{enumerate}
    \end{multicols}

    \item 敢于直面风险、战胜风险，是中华民族和中国人民的强大基因，是（ ）的精华。
    \begin{multicols}{4}
    \begin{enumerate}
        \item 中华文化
        \item 中国精神
        \item 中国品格
        \item 中华民族
    \end{enumerate}
    \end{multicols}

    \item 我们要打赢防范化解重大风险攻坚战，必须坚持和完善（ ）、推进（ ），运用制度威力应对风险挑战的冲击。
    \begin{multicols}{2}
    \begin{enumerate}
        \item 中国特色社会主义道路
        \item 中国特色社会主义制度
        \item 国家治理体系和治理能力现代化
        \item 安全和发展一体化
    \end{enumerate}
    \end{multicols}

    \item 要加强（ ）建设，坚决捍卫领土主权和海洋权益。
    \begin{multicols}{4}
    \begin{enumerate}
        \item 边防
        \item 海防
        \item 空防
        \item 陆防
    \end{enumerate}
    \end{multicols}

    \item 政治安全的核心是（ ）。
    \begin{multicols}{4}
    \begin{enumerate}
        \item 政权安全
        \item 制度安全
        \item 粮食安全
        \item 人民安全
    \end{enumerate}
    \end{multicols}

    \item 2018年4月，习近平进一步阐述了总体国家安全观，提出坚持（ ）至上的有机统一，坚持维护和塑造国家安全等重大判断。
    \begin{multicols}{2}
    \begin{enumerate}
        \item 人民安全
        \item 政治安全
        \item 国家利益
        \item 经济发展
    \end{enumerate}
    \end{multicols}
\end{enumerate}

\begin{answerbox}
多选题参考答案：
1.AB \quad 2.ABCD \quad 3.BCD \quad 4.ABCD \quad 5.AB \quad 6.BC \quad 7.ABC \quad 8.AB \quad 9.ABC
\end{answerbox}

\subsection{二、简答题}
\textbf{1. 什么是国家安全？} \\
答：国家安全是指国家政权、主权、统一和领土完整、人民福祉、经济社会可持续发展和国家其他重大利益相对处于没有危险和不受内外威胁的状态，以及保障持续安全状态的能力。

\textbf{2. 总体国家安全观的重要意义是什么？}
\begin{itemize}
    \item （1）理论意义：总体国家安全观深化了我们党对中国特色社会主义建设规律的认识，为发展马克思主义国家安全理论作出了重大原创性贡献。
    \item （2）实践意义：总体国家安全观经过实践检验、富有实践伟力，为维护和塑造新时代国家安全提供了行动纲领。
    \item （3）文化意义：总体国家安全观提炼中华优秀传统战略文化，总结我们党维护国家安全的理论和实践成果，为坚持把马克思主义基本原理同中国具体实际相结合、同中华优秀传统文化相结合树立了光辉典范。
    \item （4）世界意义：总体国家安全观开辟了国家安全治理新路径，为推动和完善全球安全治理贡献了中国方案。
\end{itemize}

\textbf{3. 中国特色国家安全道路的重要特征是什么？}
\begin{itemize}
    \item （1）坚持党的绝对领导，完善集中统一、高效权威的国家安全工作体制，实现人民安全、政治安全、国家利益至上相统一。
    \item （2）坚持捍卫国家主权和领土完整，维护边疆、边境、周边安定有序。
    \item （3）坚持安全发展，推动高质量发展和高水平安全动态平衡。
    \item （4）坚持总体战，统筹传统安全和非传统安全。
    \item （5）坚持走和平发展道路，促进自身安全和共同安全相协调。
\end{itemize}

\newpage

% ==========================================================================================
% 第十四讲
% ==========================================================================================
\section{第十四讲 坚持“一国两制”和推进祖国统一 [cite: 6]}

\subsection{重要知识点}
\begin{enumerate}
    \item “一国两制”是中国特色社会主义的一个伟大创举
    \item “一国两制”是保持港澳长期繁荣稳定的最佳制度
    \item 新时代党解决台湾问题的总体方略
\end{enumerate}

\subsection{一、选择题}

\textbf{（一）单选题}
\begin{enumerate}
    \item （ ）是两岸关系的政治基础。
    \begin{multicols}{2}
    \begin{enumerate}
        \item 继续推进两岸“三通”
        \item 促进两岸人文交流
        \item 协商解决两岸同胞关心的问题
        \item 一个中国原则
    \end{enumerate}
    \end{multicols}

    \item （ ），邓小平会见美国华人协会主席李耀滋时说：“九条方针是以叶剑英副主席的名义提出来的，实际上就是一个国家，两种制度。”这是邓小平首次提出“一国两制”的概念。
    \begin{multicols}{2}
    \begin{enumerate}
        \item 1982年1月11日
        \item 1983年2月11日
        \item 1992年1月21日
        \item 1997年2月11日
    \end{enumerate}
    \end{multicols}

    \item “一国两制”是一个完整的概念。“一国”是实行“两制”的（ ），“两制”从属和派生于“一国”并统一于“一国”之内。
    \begin{multicols}{4}
    \begin{enumerate}
        \item 目标和原则
        \item 前提和基础
        \item 方向和路径
        \item 结果和目标
    \end{enumerate}
    \end{multicols}

    \item 中央政府对特别行政区拥有（ ），这是特别行政区高度自治权的源头。
    \begin{multicols}{4}
    \begin{enumerate}
        \item 自主选择权
        \item 全面管治权
        \item 高度管理权
        \item 立法管治权
    \end{enumerate}
    \end{multicols}

    \item （ ），中共中央、国务院印发的《横琴粤澳深度合作区建设总体方案》发布，为新形势下粤澳合作开发横琴按下快进键。
    \begin{multicols}{2}
    \begin{enumerate}
        \item 2020年10月1日
        \item 2020年9月5日
        \item 2021年9月5日
        \item 2021年10月1日
    \end{enumerate}
    \end{multicols}
\end{enumerate}

\begin{answerbox}
单选题参考答案：1.D \quad 2.A \quad 3.B \quad 4.B \quad 5.C
\end{answerbox}

\textbf{（二）多选题}
\begin{enumerate}
    \item 全面准确理解和贯彻“一国两制”方针政策，必须准确把握“一国两制”的含义。“一国两制”是一个完整的概念，“一国”与“两制”的关系是（ ）。
    \begin{enumerate}
        \item “一国”是实行“两制”的前提
        \item “两制”是实行“一国”的基础
        \item “两制”从属和派生于“一国”
        \item “两制”统一于“一国之内”
    \end{enumerate}

    \item 我们要全面准确贯彻（ ）的方针，落实中央对香港、澳门特别行政区全面管治权。
    \begin{multicols}{4}
    \begin{enumerate}
        \item “一国两制”
        \item “港人治港”
        \item “澳人治澳”
        \item 高度自治
    \end{enumerate}
    \end{multicols}

    \item 特别行政区坚持实行行政主导体制，（ ）机关依照基本法和相关法律履行职责，行政机关和立法机关既互相制衡又互相配合，司法机关依法独立行使审判权。
    \begin{multicols}{4}
    \begin{enumerate}
        \item 行政
        \item 立法
        \item 司法
        \item 执法
    \end{enumerate}
    \end{multicols}

    \item “一国两制”重要制度，立足中国国情，顺应时代潮流，观照人民福祉，把、、统一起来，凝结了中国共产党人为解决国家统一问题展现出的超凡勇气和卓越智慧，是人类政治文明史上前无古人的伟大创举。
    \begin{multicols}{2}
    \begin{enumerate}
        \item 原则性和灵活性
        \item 现实性和长远性
        \item 一致性和差异性
        \item 理论性和实践性
    \end{enumerate}
    \end{multicols}

    \item “一国两制”构想在实践中首先被运用于解决香港、澳门问题。按照这一伟大构想，和香港、澳门相继回归祖国，改变了历史上但凡收复失地都要大动干戈的所谓定势，创造了推进祖国和平统一的伟大创举。
    \begin{multicols}{4}
    \begin{enumerate}
        \item 1997年
        \item 1998年
        \item 1999年
        \item 2000年
    \end{enumerate}
    \end{multicols}
\end{enumerate}

\begin{answerbox}
多选题参考答案：1.ACD \quad 2.ABCD \quad 3.ABC \quad 4.ABC \quad 5.AC
\end{answerbox}

\newpage

% ==========================================================================================
% 第十五讲
% ==========================================================================================
\section{第十五讲 推动构建人类命运共同体 [cite: 2, 3]}

\subsection{重要知识点}
\begin{enumerate}
    \item 把握世界历史大势的方法论
    \item “一国两制”是保持港澳长期繁荣稳定的最佳制度
    \item 新时代党解决台湾问题的总体方略
\end{enumerate}

\subsection{一、选择题}

\textbf{（一）单选题}
\begin{enumerate}[start=4]
    \item （ ）是构建人类命运共同体的重要思想基础，凝聚了人类不同文明的价值共识。
    \begin{multicols}{2}
    \begin{enumerate}
        \item 人类理想追求
        \item 全人类共同价值
        \item 全球经济发展
        \item 世界一体化建设
    \end{enumerate}
    \end{multicols}

    \item 新型国际关系“新”在（ ）
    \begin{multicols}{4}
    \begin{enumerate}
        \item 经济发展
        \item 合作共赢
        \item 开放包容
        \item 人类文明
    \end{enumerate}
    \end{multicols}

    \item 坚持以（ ）为依托打造全球伙伴关系。
    \begin{multicols}{2}
    \begin{enumerate}
        \item 深化外交布局
        \item 维护世界和平、促进共同发展
        \item 共商共建共享
        \item 维护党中央权威
    \end{enumerate}
    \end{multicols}

    \item 坚持以（ ）为理念引领全球治理体系变革。
    \begin{multicols}{2}
    \begin{enumerate}
        \item 深化外交布局
        \item 维护世界和平、促进共同发展
        \item 公平正义
        \item 国家核心利益
    \end{enumerate}
    \end{multicols}

    \item 坚持以（ ）为底线维护国家主权、安全、发展利益。
    \begin{multicols}{2}
    \begin{enumerate}
        \item 深化外交布局
        \item 维护世界和平、促进共同发展
        \item 公平正义
        \item 国家核心利益
    \end{enumerate}
    \end{multicols}

    \item 中国独立自主的和平外交政策，就是把（ ）放在第一位。
    \begin{multicols}{2}
    \begin{enumerate}
        \item 我国人民的利益
        \item 世界人民的利益
        \item 国家主权和安全
        \item 事情本身的是非曲直
    \end{enumerate}
    \end{multicols}

    \item 当今世界格局发展的趋势是（ ）
    \begin{multicols}{4}
    \begin{enumerate}
        \item 一极独霸世界
        \item 两极格局对峙
        \item 单边主义抬头
        \item 多极化
    \end{enumerate}
    \end{multicols}

    \item 要充分估计国际格局发展演变的复杂性，但更要看到（ ）向前推进的态势不会改变。
    \begin{multicols}{4}
    \begin{enumerate}
        \item 文化多样化
        \item 经济全球化
        \item 社会信息化
        \item 世界多极化
    \end{enumerate}
    \end{multicols}

    \item 要优化发展伙伴关系，坚持（ ）主渠道地位，让发展成果更多惠及全体人民，为世界经济全面协调可持续增长提供新动力。
    \begin{multicols}{4}
    \begin{enumerate}
        \item 东西合作
        \item 南北合作
        \item 南南合作
        \item 地区合作
    \end{enumerate}
    \end{multicols}

    \item 我们将秉持（ ）的全球治理观，积极参与全球治理体系改革和建设。
    \begin{multicols}{4}
    \begin{enumerate}
        \item 相互尊重
        \item 公平正义
        \item 合作共赢
        \item 共商共建共享
    \end{enumerate}
    \end{multicols}
\end{enumerate}

\begin{answerbox}
单选题参考答案（注：包含缺失题目1-3的答案）：
1.B \quad 2.B \quad 3.D \quad 4.B \quad 5.B \quad 6.A \quad 7.C \quad 8.D \quad 9.C \quad 10.D \quad 11.D \quad 12.B \quad 13.D \quad 14.B \quad 15.C
\end{answerbox}

\textbf{（二）多选题}
\begin{enumerate}
    \item 当今世界正经历百年未有之大变局，这样的大变局不是一时一事、一域一国之变，是（ ）。
    \begin{multicols}{4}
    \begin{enumerate}
        \item 世界之变
        \item 时代之变
        \item 历史之变
        \item 未来之变
    \end{enumerate}
    \end{multicols}

    \item 运筹好大国关系，推动构建总体稳定、均衡发展的大国关系框架至关重要。大国之间相处，要（ ）。
    \begin{multicols}{4}
    \begin{enumerate}
        \item 不冲突
        \item 不对抗
        \item 相互尊重
        \item 合作共赢
    \end{enumerate}
    \end{multicols}

    \item 和平与发展仍然是当今时代主题，（ ）成为不可阻挡的时代潮流。
    \begin{multicols}{4}
    \begin{enumerate}
        \item 和平
        \item 发展
        \item 合作
        \item 共赢
    \end{enumerate}
    \end{multicols}

    \item ( )深入发展，全球治理体系和国际秩序变革加速推进，各国相互联系和依存日益加深。
    \begin{multicols}{4}
    \begin{enumerate}
        \item 世界多级化
        \item 经济全球化
        \item 文化多样化
        \item 社会信息化
    \end{enumerate}
    \end{multicols}

    \item 20世纪90年代以来，以（ ）为核心的现代科学技术迅猛发展
    \begin{multicols}{4}
    \begin{enumerate}
        \item 计算机技术
        \item 基因技术
        \item 信息技术
        \item 生物技术
    \end{enumerate}
    \end{multicols}

    \item 中国走和平发展道路的自信和自觉（ ）
    \begin{enumerate}
        \item 来源于中华文明的深厚渊源
        \item 来源于对现实中国发展目标条件的认知
        \item 来源于对世界发展大势的把握
        \item 来源于和平发展的大局
    \end{enumerate}

    \item 和平发展是中国基于（ ）作出的战略抉择。
    \begin{multicols}{4}
    \begin{enumerate}
        \item 自身国情
        \item 社会制度
        \item 文化传统
        \item 国际形势
    \end{enumerate}
    \end{multicols}

    \item 中国走和平发展道路（ ）
    \begin{multicols}{4}
    \begin{enumerate}
        \item 符合中国根本利益
        \item 符合周边国家利益
        \item 符合世界各国利益
        \item 符合大国利益
    \end{enumerate}
    \end{multicols}

    \item 在处理国际关系和外交关系方面，我们应坚持（ ）的外交工作布局。
    \begin{multicols}{4}
    \begin{enumerate}
        \item 大国是关键
        \item 周边是首要
        \item 发展中国家是基础
        \item 多边是舞台
    \end{enumerate}
    \end{multicols}

    \item 新型国际关系的内容是（ ）
    \begin{multicols}{4}
    \begin{enumerate}
        \item 相互尊重
        \item 公平正义
        \item 合作共赢
        \item 和平共处
    \end{enumerate}
    \end{multicols}

    \item （ ）才是人类社会和平、进步、发展的永恒主题。
    \begin{multicols}{4}
    \begin{enumerate}
        \item 和平而不是战争
        \item 合作而不是对抗
        \item 共赢而不是零和
        \item 发展而不是倒退
    \end{enumerate}
    \end{multicols}

    \item 积极推动中美关系稳定发展，同俄罗斯发展全面战略协作伙伴关系，同欧洲发展（ ）伙伴关系。
    \begin{multicols}{4}
    \begin{enumerate}
        \item 和平
        \item 增长
        \item 改革
        \item 文明
    \end{enumerate}
    \end{multicols}

    \item 人类生活在同一个地球村，各国（ ）的程度空前加深。
    \begin{multicols}{4}
    \begin{enumerate}
        \item 相互联系
        \item 相互依存
        \item 相互合作
        \item 相互促进
    \end{enumerate}
    \end{multicols}

    \item 中国的治理理念和实践受到高度赞赏和广泛认同，国际（ ）进一步提高。
    \begin{multicols}{4}
    \begin{enumerate}
        \item 影响力
        \item 感召力
        \item 塑造力
        \item 发展力
    \end{enumerate}
    \end{multicols}

    \item 构建人类命运共同体，是一个（ ）的思想体系。
    \begin{multicols}{4}
    \begin{enumerate}
        \item 概念深远
        \item 科学完整
        \item 内涵丰富
        \item 意义深远
    \end{enumerate}
    \end{multicols}

    \item 构建人类命运共同体，其核心就是建设持久和平、（ ）的世界。
    \begin{multicols}{4}
    \begin{enumerate}
        \item 清洁美丽
        \item 普遍安全
        \item 共同繁荣
        \item 开放包容
    \end{enumerate}
    \end{multicols}

    \item 建设一个持久和平的世界，根本要义在于国家之间要构建（ ）的伙伴关系。
    \begin{multicols}{4}
    \begin{enumerate}
        \item 平等相待
        \item 互商互谅
        \item 互学互鉴
        \item 相互理解
    \end{enumerate}
    \end{multicols}

    \item 在国际安全新形式下，各国应树立（ ）的新安全观。
    \begin{multicols}{4}
    \begin{enumerate}
        \item 共同
        \item 综合
        \item 合作
        \item 可持续
    \end{enumerate}
    \end{multicols}

    \item 构建人类命运共同体，经济上要同舟共济，促进贸易和投资自由化便利化，推进经济全球化朝着更加（ ）共赢的方向发展。
    \begin{multicols}{4}
    \begin{enumerate}
        \item 开放
        \item 包容
        \item 普惠
        \item 平衡
    \end{enumerate}
    \end{multicols}

    \item 要维护世界贸易组织规则，支持以世界贸易组织为核心的（ ）的多边贸易体制。
    \begin{multicols}{4}
    \begin{enumerate}
        \item 开放
        \item 透明
        \item 包容
        \item 非歧视
    \end{enumerate}
    \end{multicols}
\end{enumerate}

\begin{answerbox}
多选题参考答案：
1.ABC \quad 2.ABCD \quad 3.ABCD \quad 4.ABCD \quad 5.CD \quad 6.ABC \quad 7.ABC \quad 8.ABC \quad 9.ABCD \quad 10.ABC \quad 11.ABC \quad 12.ABCD \quad 13.ABCD \quad 14.ABC \quad 15.BCD \quad 16.ABCD \quad 17.ABC \quad 18.ABCD \quad 19.ABCD \quad 20.ABCD
\end{answerbox}

\subsection{二、简答题}
\textbf{2. 全人类共同价值与普世价值的本质区别？} \\
答：（1）形成条件不同：全人类共同价值是全人类共同追求和坚守的价值观念，是人类历史从民族的历史走向世界的历史的必然产物。西方“普世价值”的本质是西方资本主义的价值观，是基于西方独特历史实践而形成的价值理念。 \\
（2）目标追求不同：全人类共同价值倡导不同社会制度、不同意识形态、不同历史文化、不同发展水平的国家在国际事务中利益共生、权利共享、责任共担。西方“普世价值”旨在推销资本主义政治体制和意识形态，所追求的是维护西方主导的国际秩序，并将之拓展到世界每一个角落。

\newpage

% ==========================================================================================
% 第十六讲
% ==========================================================================================
\section{第十六讲 全面从严治党 [cite: 9, 10]}

\subsection{重要知识点}
\begin{enumerate}
    \item 为什么要全面从严治党
    \item “四大考验”“四种危险”
    \item 如何理解全面从严治党这场伟大自我革命
\end{enumerate}

\subsection{一、选择题}

\textbf{（一）单选题}
\begin{enumerate}
    \item （ ）是新时代党的建设的根本方针。
    \begin{enumerate}
        \item 坚持党要管党
        \item 坚持党要管党、全面从严治党
        \item 全面从严治党
        \item 加强党的长期执政能力建设
    \end{enumerate}

    \item 全面从严治党以其丰富内涵诠释了（ ）的内在要求。
    \begin{multicols}{4}
    \begin{enumerate}
        \item 自我发展
        \item 自我创新
        \item 自我革命
        \item 自我改革
    \end{enumerate}
    \end{multicols}

    \item （ ）是党的政治建设的首要任务。
    \begin{enumerate}
        \item 保证全党服从中央，坚持党中央权威和集中统一领导
        \item 保证全党服从中央
        \item 坚持党的中央权威和集中统一领导
        \item 保证全党的高度团结统一
    \end{enumerate}

    \item 经过百年奋斗特别是党的十八大以来新的实践，我们党又给出了第二个答案，这就是（ ）。
    \begin{multicols}{4}
    \begin{enumerate}
        \item 反腐败斗争
        \item 自我革命
        \item 群众路线
        \item 人民监督
    \end{enumerate}
    \end{multicols}

    \item 坚持党要管党、全面从严治党的方针的关键是（ ）
    \begin{multicols}{4}
    \begin{enumerate}
        \item 全面
        \item 严
        \item 管
        \item 治
    \end{enumerate}
    \end{multicols}

    \item 党的十九大指出，党的建设要把党的（ ）摆在首位。
    \begin{multicols}{4}
    \begin{enumerate}
        \item 政治建设
        \item 思想建设
        \item 组织建设
        \item 作风建设
    \end{enumerate}
    \end{multicols}

    \item （ ）是党的基础性建设。
    \begin{multicols}{4}
    \begin{enumerate}
        \item 思想建设
        \item 组织建设
        \item 政治建设
        \item 作风建设
    \end{enumerate}
    \end{multicols}

    \item 党的作风建设的核心是：（ ）
    \begin{multicols}{2}
    \begin{enumerate}
        \item 党风问题
        \item 纪律严明
        \item 不能让享乐主义和奢靡之风卷土重来
        \item 保持党同人民群众的血肉联系
    \end{enumerate}
    \end{multicols}

    \item 我们党的最大政治优势是：（ ）
    \begin{multicols}{4}
    \begin{enumerate}
        \item 思想建设
        \item 政治建设
        \item 密切联系群众
        \item 不浮躁、不浮夸
    \end{enumerate}
    \end{multicols}

    \item （ ）审议通过《关于新形势下党内政治生活的若干准则》和《中国共产党党内监督条例》，紧紧围绕全面从严治党这个主题，继承和发扬党的优良传统和宝贵经验。
    \begin{multicols}{2}
    \begin{enumerate}
        \item 党的十八届六中全会
        \item 党的十九大
        \item 党的十九届六中全会
        \item 党的二十大
    \end{enumerate}
    \end{multicols}
\end{enumerate}

\begin{answerbox}
单选题参考答案：
1.B \quad 2.C \quad 3.A \quad 4.B \quad 5.B \quad 6.A \quad 7.A \quad 8.D \quad 9.C \quad 10.A
\end{answerbox}

\textbf{（二）多选题}
\begin{enumerate}
    \item 办好中国的事情，关键在党，关键在（ ）。
    \begin{multicols}{2}
    \begin{enumerate}
        \item 全面依法治党
        \item 人民群众监督
        \item 坚持党要管党
        \item 全面从严治党
    \end{enumerate}
    \end{multicols}

    \item 主线是加强党的长期（ ）。
    \begin{multicols}{2}
    \begin{enumerate}
        \item 执政能力建设
        \item 先进性建设
        \item 纯洁性建设
        \item 制度性建设
    \end{enumerate}
    \end{multicols}

    \item 坚持以（ ），是我们党永葆先进性、纯洁性的根本保证。
    \begin{multicols}{2}
    \begin{enumerate}
        \item 科学理论引领
        \item 用科学理论武装
        \item 以行动为准则
        \item 以人民幸福为目标
    \end{enumerate}
    \end{multicols}

    \item 党的组织建设主要包括（ ）等内容。
    \begin{multicols}{2}
    \begin{enumerate}
        \item 民主集中制建设
        \item 党的基层组织建设
        \item 干部队伍建设
        \item 党员队伍建设
    \end{enumerate}
    \end{multicols}

    \item （ ）是关系党生死存亡的问题。
    \begin{multicols}{2}
    \begin{enumerate}
        \item 党风问题
        \item 党同人民群众联系问题
        \item 党的自身发展问题
        \item 党性问题
    \end{enumerate}
    \end{multicols}

    \item 我们党在一个幅员辽阔、人口众多的发展中大国长期执政，如果没有铁的纪律，就没有党的团结统一，党的（ ）就会大大削弱，党的领导能力和执政能力就会大大削弱。
    \begin{multicols}{4}
    \begin{enumerate}
        \item 向心力
        \item 凝聚力
        \item 战斗力
        \item 组织力
    \end{enumerate}
    \end{multicols}

    \item 党的十八大以来，我们坚持党要管党、全面从严治党，坚持问题导向，以整治“四风”（ ）为突破口，大力弘扬真抓实干作风，交出了新答卷，取得了新成效。
    \begin{multicols}{4}
    \begin{enumerate}
        \item 形式主义
        \item 官僚主义
        \item 享乐主义
        \item 奢靡之风
    \end{enumerate}
    \end{multicols}

    \item 在坚持党中央权威和集中统一领导这个重大原则问题上，全党全国必须保持高度的（ ），丝毫不能动摇。
    \begin{multicols}{4}
    \begin{enumerate}
        \item 思想自觉
        \item 政治自觉
        \item 行动自觉
        \item 文化自觉
    \end{enumerate}
    \end{multicols}

    \item 新形势下党面临“四大考验”是（ ）
    \begin{multicols}{4}
    \begin{enumerate}
        \item 长期执政考验
        \item 改革开放考验
        \item 市场经济考验
        \item 外部环境考验
    \end{enumerate}
    \end{multicols}

    \item 坚持（ ），是进行具有许多新的历史特点的伟大斗争、推进中国特色社会主义伟大事业、实现民族复兴伟大梦想的根本保证。
    \begin{multicols}{4}
    \begin{enumerate}
        \item 党的领导
        \item 党要管党
        \item 全面从严治党
        \item 党的性质宗旨
    \end{enumerate}
    \end{multicols}
\end{enumerate}

\begin{answerbox}
多选题参考答案：
1.CD \quad 2.ABC \quad 3.AB \quad 4.ABCD \quad 5.AB \quad 6.BC \quad 7.ABCD \quad 8.ABC \quad 9.ABCD \quad 10.ABC
\end{answerbox}

\subsection{二、简答题}
\textbf{1. 党的政治建设的基本内容是什么？}
\begin{itemize}
    \item （1）保证全党服从中央，坚持党中央权威和集中统一领导。
    \item （2）坚定执行党的政治路线，严格遵守政治纪律和政治规矩。
    \item （3）尊崇党章，严格执行新形势下党内政治生活若干准则。
    \item （4）完善和落实民主集中制的各项制度。
    \item （5）弘扬忠诚老实、公道正派、实事求是、清正廉洁等价值观。
    \item （6）加强党性锻炼，不断提高政治觉悟和政治能力。
\end{itemize}

\end{document}